\begin{claim}{C.77}{Proved}{All abelian base strips and the exact remaining character action}
Keep the actual closed disk, proper one-dimensional
operator group $Q$, loop, principal coordinate, and
common base fixed set $A$ of C.76. For a component
$I$ of $\mathbb R\setminus A$, let
\[
        B_I=\{\beta_g|_I:g\in G\}
                    \leq\operatorname{Homeo}^+(I)
                                                               \tag{C.77.1}
\]
be its actual effective base image, with its induced
compact-open and inverse topology. Retain every strip
whose finite-endpoint characters are all zero, and
every strip for which $B_I$ is abelian. Let $Z_{\mathrm{ab}}$
be their union with all lines over $A$ and the original
circle. The restriction of the original loop to this
closed set has a based contraction through its actual
homeomorphisms, fixing the original circle and commuting
with $Q$ on its interior part. The inverse family is
jointly continuous. One compact family of original
closed-disk maps controls all evaluations, including
arbitrary accumulation of strips.

For an abelian base image there are an actual increasing
homeomorphism $b_I:I\to\mathbb R$ and a continuous
source character $d_I:G\to\mathbb R$ such that
\[
 b_I(\beta_g(t))=b_I(t)+d_I(g),\qquad
                         d_I(G)=\mathbb R.
                                                               \tag{C.77.2}
\]
The prescribed base contraction is
$b_I^{-1}(b_I(t)+r d_I(f_\tau))$, and each finite endpoint
has the simultaneous displacement $r u_\tau(a)$.
This includes abelian base strips with full-real
endpoint-stabilizer character images. The exact
restriction homomorphism is the map (C.76.4) with
$Z_{\mathrm{ab}}$ in place of $Z$; no lifting through
excluded strips is asserted. The set $Z$ of C.76 is
contained in $Z_{\mathrm{ab}}$ by the actual translation
comparison (C.75.4).

There is also an exact description of the remaining
character comparison. On a strip with nonabelian $B_I$,
write $c:G\to V\subset\mathbb R^k$ for the vector of its
finite-endpoint characters, $0\leq k\leq2$. Then
\[
 c(G_t)=V\quad(t\in I),
 \qquad
 T(g)(z,t)=(z+c(g),\beta_g(t))
                                                               \tag{C.77.3}
\]
defines a transitive action on the actual space $V\times I$.
Let $L$ be the actual closed relative group used in C.76,
and let $\mathscr P(L)$ be its Polish based path group.
For each $t\in I$, endpoint evaluation gives an open
surjection and the indicated homeomorphism
\[
 \begin{aligned}
 e_t:\mathscr P(L)&\longrightarrow V\times I,\\
 e_t(\sigma)&=(c(\sigma(1)),\beta_{\sigma(1)}(t)),\\
 \mathscr P(L)/J_t&\xrightarrow{\ \cong\ }V\times I,\\
 \sigma J_t&\longmapsto e_t(\sigma).
 \end{aligned}                                                 \tag{C.77.4}
\]
where $J_t=e_t^{-1}(0,t)$ is the actual stabilizer.
Every sufficiently small compact target set has a compact
preimage in any specified open path neighborhood whose
image contains it. No continuous section, nor uniform
compact preimage for targets escaping towards a strip
endpoint, is inferred.
\end{claim}

\begin{proof}
\emph{The actual topology of an abelian base image.}
C.75 proves that $B_I$ is transitive on $I$. Take its
closure $\overline B_I$ in the actual homeomorphism group
of $I$. It is abelian and still transitive. Its action
is free: an element fixing one point commutes with the
transitive original image and consequently fixes every
point. The verified Polish orbit-open-map theorem used
in C.30 identifies evaluation
$\overline B_I\to I$ homeomorphically, since its
stabilizer is trivial. The closure is locally compact.
The one-dimensional group comparison in (C.37.9) now
identifies its actual action, through this evaluation
homeomorphism, with the translation action of $\mathbb R$.
This is exactly C.37's transitive abelian line argument,
including its faithful action and topology hypotheses.
Because $B_I$ was already transitive, freeness of the
closure shows $B_I=\overline B_I$: the original element
with a specified evaluation value is the only element
of the closure with that value.

Choose the orientation of its real translation coordinate
so that $b_I$ is increasing. Then
\[
              d_I(g)=b_I(\beta_g(t_0))-b_I(t_0)
                                                               \tag{C.77.5}
\]
for any fixed $t_0\in I$ gives (C.77.2).
It is a character, is independent of the chosen point
$t_0$, and is continuous on the actual source group
by the displayed evaluation formula. Its kernel is
exactly $\{g:\beta_g|_I=\id\}$. Its image is all of
$\mathbb R$, by transitivity. Thus the continuous source
map and the topological isomorphism of its effective
base image have both been specified; a quotient-topology
claim for the original source restriction is unnecessary.

\emph{Three actual characters and their compact representatives.}
For an abelian base strip put
\[
       a_I(g)=(d_I(g),c_{a_1}(g),\ldots,c_{a_k}(g))
                    \in\mathbb R^{k+1},\qquad k\leq2,
                                                               \tag{C.77.6}
\]
including every finite endpoint, and write $V_I=a_I(G)$.
The three-dimensional additive path identity (C.50.3)
applies after adding zero coordinates when needed.
For every actual finite-word path from $1$ to $g$, it
places each real fraction of $a_I(g)$ in $V_I$ using
at most five signed path values. Hence $V_I$ is a real
vector subspace of dimension at most three, before
passing to any closure. It is the real span of the
original character path and is closed in this Euclidean
space. With $P=\{f_\tau\}\cup\{f_\tau^{-1}\}$ one has
\[
             r a_I(f_\tau)\in a_I(P^5)
                               \quad(0\leq r,\tau\leq1).
                                                               \tag{C.77.7}
\]

Retain the actual short-increment family $J_\delta$
of (C.76.5), and now put $E_\delta=J_\delta^{15}$.
Telescoping the original character path proves that
$a_I(J_\delta)$ spans $V_I$ for every $\delta>0$.
It is compact, symmetric and path connected.
Choose a basis of at most three vectors in this set.
The identity (C.50.3) puts the full segment from minus
to plus each basis vector in five signed values from
that set; symmetry absorbs the signs. Adding the at
most three segments therefore proves
\[
      a_I(E_\delta)\text{ contains a neighborhood of }0
                         \text{ in the actual }V_I.
                                                               \tag{C.77.8}
\]
Identity factors pad shorter products. The numbers
five and fifteen do not depend on the strip, the
dimension of its image, or the parameter values.
All original maps and inverses in $E_\delta$ tend
uniformly to the identity as $\delta\downarrow0$.

The character $d_I$ extends continuously to $L$ by
(C.77.5). The full translation equation is a closed
pointwise relation and still holds there. Each endpoint
character extends by its actual line evaluation.
Density and closedness of $V_I$ show $a_I(L)=V_I$,
with no replacement of the original image by a larger
vector space. Construct the one compact path-connected
global family $K$ from
\[
           K_0=P^5,\qquad
           K_n=E_{\delta_n}\cdots E_{\delta_1}P^5
                                                               \tag{C.77.9}
\]
as in (C.76.9). Complete compatible metrics, summable
left-product errors, compact parameter products, and
inverse control apply verbatim to these actual finite
families. Every element of $K$ is a homeomorphism of
the original closed disk fixing its circle and commuting
with the original operators.

\emph{The selection, every endpoint, and the original circle.}
On an abelian base strip use the base prescription
of the claim. Define the allowed fiber intervals by
\[
 A_n(r,\tau,t)=\operatorname{conv}
 \{u_k(t): k\in K_n,\ a_I(k)=r a_I(f_\tau)\}.
                                                               \tag{C.77.10}
\]
Equation (C.77.7) proves nonemptiness. Every label in
this set has exactly the prescribed base map by its
$d_I$ coordinate, and exactly the prescribed displacement
at every finite endpoint by the other coordinates.
For nearby parameters, the needed correction is
$r'a_I(f_{\tau'})-a_I(k)\in V_I$.
Equation (C.77.8), with an arbitrarily smaller
$\delta'\leq\delta_{n+1}$, realizes that correction
by an actual arbitrarily globally small increment.
Thus the complete lower-limit proof (C.76.14)--(C.76.15)
holds with $a_I$ in place of $c$. It uses no selected
continuous inverse of $a_I$.

On a retained nonabelian strip, all finite-endpoint
characters are zero by its definition. Use the median
base prescription and (C.76.10). The original single
short increments still lie in the new fifteenfold
family, so its lower-limit proof is unchanged.
The four parameter-face values (C.76.13) belong to
every defining interval. The relative partition-of-unity
approximation and summable limiting selection of C.76
therefore give a continuous fiber displacement $U$.
It lies in the compact limiting interval and has those
four faces exactly. Use the same explicit homeomorphism
and inverse formulas (C.76.18).

Every actual component operator still has the form
$q(s,t)=(\lambda s+\kappa,t+d)$ of C.75.
Transport the chosen abelian base coordinate by
$b_{I+d}(t+d)=b_I(t)$. This makes the transported
base character equal to $d_I$, while each transported
endpoint character is $\lambda c_a$.
Equations (C.77.2) and (C.77.10) therefore commute
with the actual $q$, including negative $\lambda$.
For proper intervals choose orbit representatives and
transport the fiber selections by multiplication by
$\lambda$, as in C.76. No nonzero component power
stabilizes a proper interval. If $I=\mathbb R$ and
the base image is abelian, the actual base translation
$t\mapsto t+d$ commutes with every $\beta_g$.
In the $b_I$ coordinate it therefore commutes with
all real translations and is itself a translation:
$b_I(t+d)=b_I(t)+e$ for one $e>0$.
The fractional base prescription consequently satisfies
$B_{r,\tau}(t+d)=B_{r,\tau}(t)+d$.
The vector $a_I$ then has only its $d_I$ coordinate,
independent of the evaluation point. Its intervals
(C.77.10) obey the exact multiplier equation
$A_n(r,\tau,t+d)=\lambda A_n(r,\tau,t)$.
Perform the selection in the actual line bundle
(C.63.20). This retains the claimed fractional base
prescription and every component equation, including
negative $\lambda$. Thus every interval and every
original operator label are covered.

At a finite endpoint $a$, every label defining an
abelian-strip fiber interval has
$u_k(a)=r u_\tau(a)$; on a zero-character strip the
same equality has both sides zero. These are exactly
the hypotheses used in (C.76.20)--(C.76.21).
Compactness of the single global $K$ gives the same
uniform endpoint estimate over varying intervals.
Every base prescription lies between the identity
and the original base map; their inverses have the
corresponding inverse bounds. Consequently the
homeomorphisms and inverses glue to (C.76.2) on all
lines over $A$, including accumulation points which
are not endpoints of any one fixed interval.

Each selected fiber value is between two $K$-values
at the same real base level. The labelled triangle
comparison therefore gives the same bounds
\[
         K_{r,\tau}(x)\in K K^{-1}K(x),\qquad
         K_{r,\tau}^{-1}(x)\in K^{-1}K K^{-1}(x).
                                                               \tag{C.77.11}
\]
On the lines over $A$, these follow from the original
prefix path as in C.76. Compactness of these actual
original families, all fixing the original circle,
proves joint boundary continuity and inverse continuity
uniformly over all strips. This verifies the claimed
contraction and the continuous restriction map in the
actual topology, by the final argument of C.76.

Every interval retained by (C.76.1) has an actual base
image conjugate to all real translations by (C.75.4)
and the nonzero character's full real image. Its base
image is therefore abelian, proving the specified
inclusion $Z\subset Z_{\mathrm{ab}}$. On the subgroup
of $\operatorname{Homeo}(Z_{\mathrm{ab}})$ preserving
$Z$, restriction to $Z$ is a continuous homomorphism
for maps and inverses. Every constructed map belongs
to that subgroup, since it preserves each strip and
each line over $A$. This actual restriction sends
the new filling to a filling of the same original
restricted loop. No identity between the two
independently selected interior homotopies is
required or asserted.

\emph{The nonabelian base and the actual transitive kernel.}
Now let $B_I$ be nonabelian, and put $D=[G,G]$.
This is a path-connected subgroup of the actual
source group: every $g,h\in G$ has the finite-word
paths (C.75.7), so their commutator has a commutator
path from the identity; finite products give paths
to every element of $D$. Its base image is nontrivial,
normal in $B_I$, and path connected. Its common
fixed set is $B_I$-invariant, and hence is either
empty or all of $I$, since $B_I$ is transitive.
Nontriviality excludes the latter.

For every point of $I$, an actual commutator-product
path moves it to a different point; its inverse
path supplies movement to the other side. The
evaluation interval therefore contains a neighborhood
of that point. Each $\beta(D)$-orbit is path connected
and open. These disjoint orbits partition the
connected interval $I$, so there is exactly one.
Thus $\beta(D)$ is transitive. All endpoint characters
vanish on $D$, because they are additive characters.

For $v\in V=c(G)$ choose an actual $g\in G$ with
$c(g)=v$. For any desired $t'\in I$, choose $h\in D$
with $\beta_h(\beta_g(t))=t'$. Then $hg$ has character
$v$ and the desired base value. Taking $t'=t$ proves
$c(G_t)=V$, and arbitrary $t'$ proves (C.77.3).
The inverse of its displayed map is exactly
$(z,t)\mapsto(z-c(g),\beta_g^{-1}(t))$.
Both map and inverse depend continuously on the
actual source element on every compact target set,
by continuity of the character and of base restriction.
Its action kernel is exactly
\[
             \{g\in G:c(g)=0, \beta_g|_I=\id\}.
                                                               \tag{C.77.12}
\]
This proves the statement for the actual source words,
not only for their closure.

The vector space $V$ and the endpoint characters
extend to $L$ as in C.76. Endpoint action gives a
continuous action of the Polish based path group
$\mathscr P(L)$ on $V\times I$. Every original word
$hg$ just used is reached by its finite-word path,
so this path-group action is transitive. The target
is the nonmeagre Hausdorff manifold $V\times I$.
The exact orbit-open-map theorem already verified
in C.30 makes $e_t$ in (C.77.4) open. Its fibers
are the actual cosets of its closed stabilizer
$J_t$. The displayed quotient map is consequently
a continuous open bijection and hence a homeomorphism.
The kernel of the whole path-group action consists
of the paths whose endpoints satisfy (C.77.12),
now read in $L$; it has not been suppressed.

Apply C.64's proved compact-covering lemma to this
open map, using the complete compatible path metric.
For every compact target set lying in the image of
an assigned open path neighborhood, it provides
a compact preimage entirely inside that neighborhood.
The union of the values of these compact paths is
also compact in the original relative group, including
the inverse maps. This proves the final compactness
assertion. The argument supplies no single compact
family for a collection of target points escaping
to a finite endpoint or to infinity in the base
coordinate. The original-boundary-controlled selection
on these remaining nonabelian strips is therefore
still a precise unresolved step in C.12.
\end{proof}
